\documentclass[11pt]{article}
\usepackage[margin=0.9in]{geometry}
\usepackage{fontspec}
\usepackage{kotex}
\usepackage{amsmath}
\usepackage{longtable}
\usepackage{array}
\usepackage{setspace}
\usepackage[hidelinks]{hyperref}
\setmainfont{Times New Roman}
\setstretch{1.05}
\emergencystretch=3em
\title{Repeated Identical In-District Early-Vote Pairs That Overwhelm the Random-Chance Hypothesis}
\author{Junho Kim \\ No affiliation \\ e-mail: junhokim0331@gmail.com}
\date{June 2026}
\begin{document}
\maketitle


\section*{Abstract}

This paper examines whether the repeated identical vote-count pairs observed in in-district early-vote counting units after the June 3, 2026 Korean local election can be explained by ordinary chance. The starting point is the public report of identical vote counts in twelve places nationwide, including ten Gwangju-Jeonnam rows where the same two leading candidates received identical counts across paired counting units. This study rechecks the twelve event rows by extracting them directly from the National Election Commission election-statistics system's official counting-unit HTML pages. The Gwangju-Jeonnam case is defined as five repeated identical vote pairs within the same election, the same top-two candidate combination, and the same in-district early-vote category.

To build a historical baseline, this study parses 81,701 early-vote counting-unit rows from official National Election Commission and public-data files from 2014 through 2025. The primary benchmark is the actual top-two candidate vote-pair distribution in governor races. Across 51 governor constituency contests and approximately 1,514,172 actual comparison pairs, only 15 identical top-two vote pairs are observed, and the maximum repeated-pair count within any one constituency contest is three. Under the expanded benchmark, the probability of observing at least five repeated pairs in the 2026 Gwangju-Jeonnam setting is approximately 0.115\% using a Poisson approximation with $N=393$ and $K=100,944.8$. In a nonparametric resampling test drawing 393 units without replacement from the historical governor top-two pair pool, no simulation among 200,000 trials produces five or more repeated pairs; the rule-of-three 95\% upper bound is approximately 0.0015\%.

This paper also checks the difference between early-vote and election-day two-party vote shares in the 2016, 2020, and 2024 National Assembly elections. In 2020 and 2024, every analyzable constituency shows a higher Democratic-party two-party share in early voting than on election day. This auxiliary test is not treated as direct proof of a specific misconduct mechanism, but it reinforces the need for a broader raw-data audit. The conclusion is not a legal finding that a specific act of misconduct has been proven. It is a narrower statistical conclusion: even under conservative assumptions, the public data contain anomalies that weaken the random-chance hypothesis and require independent access to raw counting records.

Keywords: in-district early voting, identical vote pair, election integrity, duplicate probability, Poisson approximation, early-vote deviation, raw-data audit

\section{Introduction}

After the 2026 Korean local election, several reports noted that two leading candidates received exactly the same vote counts in different in-district early-vote counting units. A TV Chosun report dated June 8, 2026 stated that identical vote counts appeared in twelve places nationwide, including ten Gwangju-Jeonnam rows where Democratic Party candidate Min Hyung-bae and People Power Party candidate Lee Jung-hyun received the same paired vote counts. The same report also presented the Incheon Songdo 1-dong and Songdo 2-dong case, where Democratic Party candidate Park Chan-dae received 3,030 votes and People Power Party candidate Yoo Jeong-bok received 1,440 votes in both units.

The subject of this paper is statistical rather than partisan. The questions are:

1. How often has a pattern of five or more identical top-two candidate vote pairs occurred in the same election, same candidate combination, and same in-district early-vote category in historical official data?
2. Under a simple duplicate-probability model, how likely is the observed repetition?
3. Can the observed pattern be closed as ordinary random coincidence, or does it constitute a statistical anomaly requiring raw-data verification?

This paper answers the third question cautiously. Statistical improbability is not a sufficient condition for proving misconduct. However, when the improbability is strong enough, an election-management institution cannot reasonably end the matter by merely saying "coincidence." It must produce the raw records that can either explain or falsify the anomaly.

\subsection{Research Hypotheses}

The primary test concerns repeated identical vote pairs.

\begin{itemize}
\item Null hypothesis $H_0$: The five identical top-two candidate vote pairs observed in the 2026 Gwangju-Jeonnam in-district early-vote results are compatible with the duplicate-pair structure observed in past governor elections. In plain terms, the event is something that historical data would make plausible.
\item Alternative hypothesis $H_1$: The five identical top-two candidate vote pairs in the 2026 Gwangju-Jeonnam in-district early-vote results are difficult to reconcile with the historical empirical distribution. In plain terms, the event is an outlier relative to past official data.
\end{itemize}

The secondary test concerns the difference between early voting and election-day voting.

\begin{itemize}
\item Null hypothesis $H_0'$: The vote-share difference between early voting and election-day voting is explainable within ordinary constituency-level voter self-selection and voting-behavior differences.
\item Alternative hypothesis $H_1'$: The early-vote and election-day vote-share difference repeats in the same direction across many constituencies and remains insufficiently explained without additional covariates or raw-data verification.
\end{itemize}

Both $H_1$ and $H_1'$ are treated as statistical audit triggers, not legal judgments.

\section{Definition of the Observed Event}

The event analyzed in this paper is defined as follows:

\begin{itemize}
\item Election: the 9th nationwide local election held in 2026.
\item Vote type: in-district early voting.
\item Comparison unit: eup, myeon, or dong counting-unit results within the same metropolitan or provincial governor race.
\item Candidate rule: the principal first- and second-place candidates in the relevant race.
\item Event: two different counting units report the same first-candidate vote count and the same second-candidate vote count simultaneously.
\end{itemize}

The reported Gwangju-Jeonnam case is summarized as five pairs. This study re-extracts and checks these values from the official NEC election-statistics counting-unit pages.

\begin{center}
\small
\begin{tabular}{r p{0.29\textwidth} p{0.29\textwidth} r r}
\hline
Pair & Counting unit A & Counting unit B & Min Hyung-bae & Lee Jung-hyun \\
\hline
1 & Songjeong 1-dong, Gwangsan-gu, Gwangju & Geumsan-myeon, Goheung-gun & 1,401 & 120 \\
2 & Haui-myeon, Sinan-gun & Samil-dong, Yeosu-si & 506 & 42 \\
3 & Iyang-myeon, Hwasun-gun & Byeongyeong-myeon, Gangjin-gun & 444 & 46 \\
4 & Eomda-myeon, Hampyeong-gun & Bukha-myeon, Jangseong-gun & 606 & 57 \\
5 & Nodong-myeon, Boseong-gun & Palgeum-myeon, Sinan-gun & 356 & 42 \\
\hline
\end{tabular}
\end{center}

The Incheon Songdo case is treated separately because it belongs to another contest.

\begin{center}
\small
\begin{tabular}{l l r r}
\hline
Counting unit A & Counting unit B & Park Chan-dae & Yoo Jeong-bok \\
\hline
Songdo 1-dong & Songdo 2-dong & 3,030 & 1,440 \\
\hline
\end{tabular}
\end{center}

The paper therefore separates two levels:

\begin{itemize}
\item Internal Gwangju-Jeonnam event: five pairs in one contest environment.
\item Nationwide reported event: six pairs across more than one contest.
\end{itemize}

The stronger statistical target is the internal Gwangju-Jeonnam event. The nationwide six-pair set is more exposed to post-search and multiple-comparison objections because it was found after scanning across contests.

\section{Data and Reproduction Method}

\subsection{Historical Baseline Data}

The study uses public National Election Commission and public-data election counting files from 2014 through 2025.

\begin{center}
\small
\begin{tabular}{l r}
\hline
Election & Parsed rows \\
\hline
2014 local & 25,600 \\
2016 National Assembly & 3,661 \\
2017 presidential & 3,684 \\
2018 local & 22,926 \\
2020 National Assembly & 3,681 \\
2022 presidential & 3,703 \\
2022 local & 10,942 \\
2024 National Assembly & 3,754 \\
2025 presidential & 3,750 \\
Total & 81,701 \\
\hline
\end{tabular}
\end{center}

Aggregate rows, subtotal rows, shipboard/residential special votes, overseas votes, and out-of-country categories are excluded from the main comparison. In-district early voting and out-of-district early voting are separated, and the central analysis uses in-district early-vote rows.

At the time this paper was prepared, the official integrated XLSX file for the 2026 local election was not available in the same public-data format used for the historical baseline. However, the NEC election-statistics system was providing counting-unit results for the 9th nationwide local election. This study saves and parses the relevant POST-response HTML from the official "counting-unit result" pages and rechecks the twelve event rows. The procedure is implemented in \path{scripts/fetch_nec_2026_duplicate_cases.py}. The main outputs are:

\begin{itemize}
\item \path{outputs/nec_2026_reported_duplicate_cases.csv}
\item \path{outputs/nec_2026_reported_duplicate_pairs.csv}
\end{itemize}

Thus the 2026 event values no longer depend only on a media screenshot. When the official integrated XLSX file or original counting statements become available, the same event definition should be re-applied.

\subsection{Inclusion and Exclusion Criteria}

The 2014 local election is included to increase baseline reliability. It is the first nationwide local election after the introduction of the nationwide early-voting system, and the public counting files distinguish in-district and out-of-district early voting. The 2010 local election and the 2012 presidential/National Assembly elections are excluded from the main baseline because they do not belong to the same in-district early-voting institutional regime.

This choice does not bias the result in favor of anomaly detection. Including the 2014 local election increases the number of historical duplicate pairs and therefore weakens the apparent rarity of the 2026 event. Even after this broader inclusion, the 2026 Gwangju-Jeonnam five-pair repetition exceeds the historical maximum.

\subsection{Actual First- and Second-Place Candidate Rule}

A naive analysis might compare the first two candidates listed in each file. That does not always match the reported meaning of "first and second place." This paper therefore uses the following procedure:

1. Group counting units by election, office type, constituency, and vote type.
2. Compute total votes for each candidate within the group.
3. Select the actual first- and second-place candidates by total votes.
4. For each eup, myeon, or dong in-district early-vote counting unit, form the ordered vote pair for these two candidates.
5. Count whether the same vote pair appears in two or more different units.

This is the rule closest to the reported event.

\subsection{Mapping Claims to Verification Files}

This paper separates statistically reproducible claims from administrative or legal claims that require additional raw records. Without this separation, a critic can attack the weakest part of a broad claim and dismiss the entire analysis. The reproducible claims in this package are:

1. The twelve 2026 event rows are reproduced from official NEC election-statistics HTML pages.
2. Those twelve rows form six identical vote pairs.
3. In the 51 historical governor in-district early-vote contests, the maximum repeated-pair count within a single contest is three.
4. With $N=393$ and $K=100,944.8$, the Poisson-tail probability for five or more repeated pairs in Gwangju-Jeonnam is approximately 0.115\%.
5. In 200,000 nonparametric trials sampling 393 units without replacement from the historical governor actual-pair pool, five or more repeated pairs occur zero times.
6. In the 2020 and 2024 National Assembly elections, the Democratic candidate's early-vote two-party share is higher than the election-day two-party share in every analyzable constituency.

The unresolved causal claim is different: whether the identical pairs were produced by natural counting processes, data-entry or display errors, review/reclassification procedures, or another cause. That question requires original counting statements, ballot-sorter first-pass outputs, review-ballot allocation records, and input logs. The detailed claim-to-file mapping is provided in \texttt{evidence\_matrix\_ko.md}.

The completed conclusion of this study is therefore not "the cause has been proven." It is: "official page values and historical baselines are sufficient to define a statistical anomaly requiring raw-data audit."

\section{Empirical Baseline}

\subsection{Historical Governor Elections}

The closest comparison class for the 2026 Gwangju-Jeonnam case is past governor elections because the observed event occurred in a metropolitan/provincial governor race and in the in-district early-vote category.

Reanalysis of the 2014, 2018, and 2022 governor-election in-district early-vote data under the actual top-two candidate rule gives:

\begin{center}
\small
\begin{tabular}{l r}
\hline
Quantity & Value \\
\hline
Governor constituency contests & 51 \\
Total comparison pairs & 1,514,172 \\
Identical vote pairs & 15 \\
Constituency contests with at least one identical pair & 9 \\
Maximum identical pairs within one contest & 3 \\
Constituency contests with five or more pairs & 0 \\
\hline
\end{tabular}
\end{center}

The observed historical identical pairs are:

\begin{center}
\scriptsize
\begin{longtable}{p{0.10\textwidth} p{0.10\textwidth} p{0.15\textwidth} p{0.15\textwidth} p{0.10\textwidth} p{0.24\textwidth}}
\hline
Election & Constituency & First-place candidate & Second-place candidate & Pair & Units \\
\hline
2014 local & Gangwon & Choi Moon-soon & Choi Heung-jip & 49, 80 & Duchon-myeon, Gimsatgat-myeon \\
2014 local & Gangwon & Choi Moon-soon & Choi Heung-jip & 57, 95 & Buron-myeon, Mukho-dong \\
2014 local & Gyeongnam & Hong Joon-pyo & Kim Kyung-soo & 293, 83 & Bugok-myeon, Ungyang-myeon \\
2014 local & Gyeongbuk & Kim Kwan-yong & Oh Joong-gi & 160, 23 & Chojeon-myeon, Hyeondong-myeon \\
2014 local & Gyeongbuk & Kim Kwan-yong & Oh Joong-gi & 186, 28 & Jisan-dong, Andeok-myeon \\
2014 local & Gyeongbuk & Kim Kwan-yong & Oh Joong-gi & 342, 43 & Pyeonghwa-dong, Bongyang-myeon \\
2014 local & Jeonnam & Lee Nak-yon & Lee Seong-su & 124, 18 & Nam-myeon, Eomda-myeon \\
2014 local & Jeonnam & Lee Nak-yon & Lee Seong-su & 190, 33 & Bugil-myeon, Heuksan-myeon \\
2014 local & Jeonbuk & Song Ha-jin & Park Cheol-gon & 145, 41 & Ungpo-myeon, Boan-myeon \\
2014 local & Jeonbuk & Song Ha-jin & Park Cheol-gon & 174, 27 & Ipyeong-myeon, Ibaek-myeon \\
2014 local & Jeonbuk & Song Ha-jin & Park Cheol-gon & 174, 50 & Seosu-myeon, Jungang-dong \\
2014 local & Chungnam & Ahn Hee-jung & Chung Jin-suk & 54, 61 & Sicho-myeon, Munsan-myeon \\
2018 local & Daegu & Kwon Young-jin & Im Dae-yoon & 674, 559 & Gwaneum-dong, Hyeonpung-myeon \\
2022 local & Jeonnam & Kim Young-rok & Lee Jung-hyun & 377, 90 & Dongbok-myeon, Jangsan-myeon \\
2022 local & Chungbuk & Kim Young-hwan & Noh Young-min & 319, 158 & Yangsan-myeon, Naebuk-myeon \\
\hline
\end{longtable}
\end{center}

The point is not that historical duplicate pairs never occur. They do. In particular, including the 2014 local election increases the historical duplicate count. But no historical governor contest in the expanded baseline has five repeated pairs within one contest. The historical maximum is three.

A predictable objection is that if three pairs occurred historically, five pairs should not be treated as special. That objection confuses a maximum observation with a tail probability. Three pairs is the upper end observed after pooling 51 historical governor contests. Five pairs exceeds that observed upper end by two additional repeated pairs. Under $N=393$ and $K=100,945$, three or more pairs has a probability near 4.23\%; five or more pairs has a probability near 0.115\%. This paper therefore does not argue that "three pairs are already proof, so five pairs are proof." It accepts three as an observed rare historical upper value and treats five as an event beyond that upper value.

\subsection{Empirical Estimate of the Effective Vote-Pair Space}

The historical governor dataset contains 1,514,172 total comparison pairs and 15 observed identical vote pairs. If the effective number of possible first-second vote pairs is denoted by $K$, then the expected number of collisions is roughly:

\[
E(C) \approx \frac{N(N-1)}{2K}.
\]

Solving backward from the historical governor data gives:

\[
\hat{K} \approx \frac{1,514,172}{15} \approx 100,945.
\]

This is not a strict structural parameter. Real election data are not generated from identical independent draws; unit size, geography, party support, and turnout matter. But $K$ is an empirical reference value reflecting how often actual Korean governor-election top-two vote pairs collided in official data.

\section{Probability Calculation}

\subsection{Poisson Approximation}

Suppose one constituency-level contest contains $N$ in-district early-vote counting units, and each unit produces a top-two vote pair from an effective pair space of size $K$. Let $C$ be the number of identical vote-pair collisions among distinct unit pairs. For rare collision events, $C$ can be approximated by:

\[
C \sim \mathrm{Poisson}(\lambda),
\]

with

\[
\lambda = \frac{N(N-1)}{2K}.
\]

The probability of interest is:

\[
P(C \geq 5),
\]

which corresponds to at least five identical vote-pair repetitions within the Gwangju-Jeonnam setting.

\subsection{Baseline with $N=393$ and $K=100,945$}

Here $N$ is not the number of already reported rows. It is the number of potential in-district early-vote counting units that could have produced identical pairs within the same contest environment. This means the eup, myeon, and dong level units in the same candidate combination and vote type.

This study does not estimate $N$ from administrative intuition. It counts $N$ directly from the NEC VCCP08 official HTML pages. After parsing all governor counting-unit pages for the five Gwangju districts and the twenty-two Jeonnam cities/counties, and counting rows that are in-district early-vote rows rather than aggregate rows, the output \texttt{outputs/nec\_2026\_gwangju\_jeonnam\_unit\_counts.csv} gives 96 Gwangju units and 297 Jeonnam units, for a total of 393. Therefore the baseline calculation uses $N=393$.

$K=100,945$ is also not arbitrary. In the historical governor baseline, there are 1,514,172 actual top-two comparison pairs and 15 observed identical vote pairs. Reversing the rare-collision formula gives $K$ near 100,945.

\[
\hat{K} \approx \frac{1,514,172}{15} \approx 100,945
\]

Thus $N=393$ and $K=100,945$ are not numbers inserted for convenience. They combine the Gwangju-Jeonnam in-district early-vote unit count measured from official NEC pages with the empirically inferred effective pair space from historical governor data.

Using these values:

\[
\lambda = \frac{393 \times 392}{2 \times 100,945} \approx 0.76307.
\]

The computed probabilities are:

\begin{center}
\small
\begin{tabular}{l r r}
\hline
Event & Probability & Reciprocal interpretation \\
\hline
At least 5 pairs & about 0.115\% & about 1 in 871 \\
At least 6 pairs & about 0.0143\% & about 1 in 6,982 \\
\hline
\end{tabular}
\end{center}

In this calculation, five pairs refers to the internal Gwangju-Jeonnam event. Six pairs is a simplified reference value for the nationwide reported event. The nationwide six-pair figure requires additional multiple-search adjustment because it was found after looking across more than one contest.

\subsection{Sensitivity to the Effective Pair Space}

The probability changes depending on how the effective pair space $K$ is set. With $N=393$, the sensitivity calculation is:

\begin{center}
\small
\begin{tabular}{r r r}
\hline
Effective pair space $K$ & $P(C \geq 5)$ & $P(C \geq 6)$ \\
\hline
50,000 & 2.0550\% & 0.5056\% \\
100,000 & 0.1197\% & 0.0151\% \\
100,945 & 0.1148\% & 0.0143\% \\
200,000 & 0.0051\% & 0.00033\% \\
337,354 & 0.00043\% & 0.000016\% \\
500,000 & 0.000064\% & 0.0000016\% \\
\hline
\end{tabular}
\end{center}

Even under a highly conservative $K=50,000$ assumption, which makes duplicate pairs much easier to produce, the probability of at least five pairs is about 2.05\%. That is not impossibility, but it remains a low-tail event. Under the expanded historical governor benchmark, $K=100,945$, the probability is about 0.115\%. This is sufficiently rare to function as an audit trigger even after including the smaller early-vote scale of 2014 in the historical baseline.

\subsection{Sensitivity to the Number of Counting Units}

The official HTML-based baseline is $N=393$. To check whether the conclusion depends narrowly on that value, the paper also holds $K=100,945$ fixed and varies $N$ from 350 to 450.

\begin{center}
\small
\begin{tabular}{r r r r}
\hline
Units $N$ & $\lambda$ & $P(C \geq 5)$ & Reciprocal interpretation \\
\hline
350 & 0.6050 & 0.0410\% & about 1 in 2,441 \\
390 & 0.7515 & 0.1074\% & about 1 in 931 \\
393 & 0.7631 & 0.1148\% & about 1 in 871 \\
400 & 0.7905 & 0.1340\% & about 1 in 746 \\
430 & 0.9137 & 0.2500\% & about 1 in 400 \\
450 & 1.0008 & 0.3672\% & about 1 in 272 \\
\hline
\end{tabular}
\end{center}

The table shows that the assumed unit count changes the exact probability but not the direction of the conclusion. Even at $N=450$, at least five repeated pairs remains below 0.4\%. The Gwangju-Jeonnam five-pair result therefore does not depend only on the single $N=393$ baseline.

\subsection{Nonparametric Resampling Test}

The Poisson approximation is useful because it is easy to interpret, but it still contains the model assumption of an effective pair space $K$. To reduce model dependence, this paper directly resamples actual historical top-two vote pairs from governor in-district early-vote data.

The procedure is:

1. Build the pool of 10,322 actual top-two candidate vote pairs from the 2014, 2018, and 2022 governor in-district early-vote data.
2. Draw 393 units without replacement.
3. Count the number of repeated vote pairs in the draw.
4. Repeat this procedure 200,000 times with a fixed random seed.

Sampling without replacement prevents artificial duplicates caused by drawing the same historical unit more than once. The test directly asks: if 393 units are drawn from actual past governor early-vote units, how often do five or more repeated pairs appear?

The observed simulation results are:

\begin{center}
\small
\begin{tabular}{l r r}
\hline
Threshold & Hits in 200,000 trials & Empirical probability \\
\hline
$C \geq 3$ & 77 & 0.000385 \\
$C \geq 4$ & 4 & 0.000020 \\
$C \geq 5$ & 0 & 0 at simulation resolution \\
\hline
\end{tabular}
\end{center}

Zero hits in 200,000 trials does not prove mathematical impossibility. It means that the event is below the simulation resolution. Applying the rule of three gives a rough 95\% upper bound of about $3/200,000 = 0.000015$, or 0.0015\%.

This result is stronger than the Poisson approximation, but the paper does not rely on it alone. The main point is convergence across approaches: the historical maximum is three pairs, the Poisson approximation gives 0.115\% for five or more pairs, and direct resampling yields zero occurrences in 200,000 trials. All three checks reject the interpretation that the Gwangju-Jeonnam five-pair cluster is an ordinary event.

\section{Why One Pair and Five Pairs Are Different}

The most common misunderstanding in identical-vote-pair discussions is a linear comparison: if one pair can occur, then several pairs can also occur. But the event tested here is not whether at least one identical pair exists. The tested event is how many identical pairs accumulate within the same contest, same candidate combination, and same in-district early-vote category. The former is a broad collision event. The latter is an upper-tail repetition event.

When the number of counting units is large, one identical pair can easily appear because the number of possible unit-to-unit comparisons grows quickly.

\[
\binom{N}{2} = \frac{N(N-1)}{2}
\]

For $N=393$, there are 77,028 unit pairs. In a comparison space that large, one matching vote pair somewhere is not by itself an anomaly. This paper therefore does not treat one identical pair as suspicious by itself.

Five pairs are a different object. If $C$ is the number of identical vote-pair collisions in the contest environment, the relevant question is not $C \ge 1$, but $C \ge 5$. In the Poisson approximation, the probability of at least $m$ repeated pairs is:

\[
P(C \ge m) = 1 - e^{-\lambda}\sum_{c=0}^{m-1}\frac{\lambda^c}{c!}
\]

As $m$ rises from 1 to 3, 4, and 5, the probability does not decline in a merely linear way. It moves rapidly into the tail.

The distinction is visible in the numbers. Under the baseline model, $P(C \geq 3)$ is about 4.23\%, while $P(C \geq 5)$ is about 0.115\%. Both are below the center of the distribution, but five is roughly 37 times rarer than three under the same assumptions. Historical data also show this difference: three pairs occurred as the historical maximum; five did not occur in any historical governor contest in the baseline.

\begin{center}
\small
\begin{tabular}{l r r p{0.27\textwidth}}
\hline
Threshold & Probability & Intuitive frequency & Interpretation \\
\hline
$C \ge 1$ & about 53.37\% & about 1 in 1.9 & one pair is plausible \\
$C \ge 2$ & about 17.88\% & about 1 in 5.6 & uncommon but unsurprising \\
$C \ge 3$ & about 4.23\% & about 1 in 23.7 & observed historical upper end \\
$C \ge 4$ & about 0.773\% & about 1 in 129 & strong audit signal \\
$C \ge 5$ & about 0.115\% & about 1 in 871 & event above the historical maximum \\
\hline
\end{tabular}
\end{center}

Thus the difference between three pairs and five pairs is not merely "two more pairs." Under the same baseline, at least five pairs is about 37 times rarer than at least three pairs. More importantly, three pairs is the rare upper end actually observed in the historical governor data, whereas five pairs exceeds that upper end.

This distinction matters for rebuttal. Saying "if three pairs occurred historically, five pairs are also fine" is like saying that because 38 degrees Celsius has occurred historically, 40 degrees must be the same kind of event. The numeric difference may look small, but exceeding the historical maximum changes the statistical interpretation. In the same way, three pairs set the observed benchmark, while five pairs test and exceed it.

Another key issue is post-search scope. If one searches all elections, all offices, all candidates, all vote types, and all regions, then a rare-looking coincidence will eventually be found. This paper therefore places the main claim on the Gwangju-Jeonnam internal five-pair cluster rather than the entire nationwide set. The Gwangju-Jeonnam cluster is defined by the same office, same top-two candidate pair, same in-district early-vote category, and the same contest environment.

In conclusion, this paper does not say that one identical pair is suspicious by itself. One pair is possible. Three pairs are also acknowledged as a rare historical upper end. The problem is that five pairs are concentrated in the same structure and exceed both the historical maximum and the tail-probability benchmark.

\section{Specificity of the Incheon Songdo Case}

The Incheon Songdo case is not simply another row to be mechanically added to the Gwangju-Jeonnam cluster. It belongs to a different contest and should be analyzed separately.

Reports state that Songdo 1-dong and Songdo 2-dong had the same final in-district early-vote result for the two leading candidates: Park Chan-dae 3,030 votes and Yoo Jeong-bok 1,440 votes. However, the electorate counts, third-candidate votes, invalid votes, and abstentions differed. Public election-management explanations also stated that the first-pass ballot-sorter results were not identical, and that the final equality emerged only after reviewed ballots were hand-counted and added.

The reported first-pass results were:

\begin{center}
\small
\begin{tabular}{l r r r}
\hline
Unit & Park first pass & Yoo first pass & Reviewed ballots \\
\hline
Songdo 1-dong & 3,016 & 1,427 & 44 \\
Songdo 2-dong & 3,011 & 1,429 & 53 \\
\hline
\end{tabular}
\end{center}

The final result was:

\begin{center}
\small
\begin{tabular}{l r r}
\hline
Unit & Park Chan-dae final & Yoo Jeong-bok final \\
\hline
Songdo 1-dong & 3,030 & 1,440 \\
Songdo 2-dong & 3,030 & 1,440 \\
\hline
\end{tabular}
\end{center}

The statistical specificity of Songdo can be viewed in two ways. First, using the same $K=100,944.8$ benchmark, Yeonsu-gu has 15 in-district early-vote units in the official HTML parsing. The number of unit comparisons is $\binom{15}{2}=105$, and:

\[
\lambda_{\text{Yeonsu}} = \frac{105}{100,944.8} \approx 0.001040
\]

Therefore the probability that at least one identical vote pair appears somewhere inside the 15 Yeonsu in-district early-vote units is:

\[
P(C \ge 1) \approx 1-e^{-0.001040} \approx 0.104\%
\]

This is about 1 in 962.

Second, if Songdo 1-dong and Songdo 2-dong are treated as the already specified pair of interest, the conditional probability that those two units have the same top-two vote pair is:

\[
\frac{1}{K} \approx \frac{1}{100,944.8} \approx 0.000991\%
\]

This is about 1 in 100,945. That number is appropriate only when the two Songdo units are treated as pre-specified. If the event is instead defined after searching nationwide, post-search adjustment is required.

The Songdo case is therefore not part of the same primary test as the Gwangju-Jeonnam five-pair cluster, but it is a strong auxiliary audit signal. The key question is not only whether final numbers can happen to match. It is how different first-pass results and reviewed-ballot pools produced the same final top-two vote pair.

The same fact can be read in two directions. From the election-management perspective, different first-pass results and reviewed-ballot counts can be offered as an explanation that the final equality was not a simple copy or display error. From an audit perspective, the same explanation makes the reviewed-ballot candidate allocation records essential.

The core verification material for Songdo is therefore not the final published table alone. It is the candidate-by-candidate allocation of the reviewed ballots.

\section{Statistical Anomaly Versus Evidence of Misconduct}

It is important to separate two propositions:

1. The repeated identical vote-pair pattern is statistically anomalous under the historical benchmark.
2. A specific person or institution committed a specific act of misconduct.

This paper supports the first proposition. It does not, by itself, prove the second. The phrase "evidence" should therefore be used in two layers. First, identical vote-pair repetition is statistical evidence against the ordinary-randomness explanation. Second, it is not direct evidence identifying a perpetrator, intent, method, or legal violation.

This distinction does not weaken the audit conclusion. It makes the conclusion harder to dismiss. A statistical anomaly need not prove a crime in order to justify raw-data disclosure. Election administration is a public trust system; when an official result produces a rare pattern under a transparent benchmark, the institution must be able to show the records that explain it.

\section{Raw Materials Required for Verification}

To transform the statistical question into an administrative or legal conclusion, the following records are needed:

1. Original counting statements for all twelve event rows.
2. Ballot-sorter first-pass results for the affected units.
3. Reviewed-ballot counts and candidate allocation records.
4. Manual recount or hand-count logs, if any.
5. Data-entry logs, correction logs, and publication/export logs.
6. Ballot-box movement and chain-of-custody records.
7. Hashable raw exports for the official integrated counting-unit files.

If these records show a coherent unit-by-unit path from first pass to final publication, the anomaly can be explained. If they do not, the statistical anomaly becomes stronger. The present paper cannot resolve that factual question without access to those raw materials.

\section{Limitations}

The study has several limitations.

First, the 2026 integrated XLSX file comparable to the historical public-data files was not available at the time of writing. The 2026 event rows are rechecked through official NEC HTML pages, but future work should repeat the analysis using the official integrated file or original counting statements.

Second, the Poisson model simplifies the data-generating process. Real voting units are not identical independent draws. Size, geography, turnout, age composition, political preference, and local campaign effects matter. That is why the paper also uses historical actual-pair resampling rather than relying only on a theoretical model.

Third, the early-vote versus election-day vote-share analysis is auxiliary. A systematic early-vote advantage for one party can arise from voter self-selection, party mobilization, demographic differences, pandemic-era voting behavior, and other lawful causes. The analysis is included because the one-directional pattern is statistically notable, not because it independently proves misconduct.

Fourth, the paper does not claim to identify a mechanism. The candidates include possibilities such as ordinary coincidence, administrative correction, display error, data-entry error, review-ballot allocation, or deliberate manipulation. Mechanism identification requires raw records.

\section{Early-Vote Versus Election-Day Vote-Share Difference}

Separate from repeated identical vote pairs, this study uses early-vote versus election-day vote-share differences as an auxiliary test. The purpose is to check whether early voting reflects the same direction and scale of political choice as election-day voting. For a given constituency, let $p_E$ be a candidate's early-vote share, $p_D$ the same candidate's election-day vote share, $n_E$ the early-vote count, and $n_D$ the election-day vote count. Under a null hypothesis that the two vote groups are random samples from the same underlying support rate $p$, the standard two-proportion z statistic is:

\[
z = \frac{p_E - p_D}{\sqrt{\hat{p}(1-\hat{p})(1/n_E + 1/n_D)}}
\]

where $\hat{p}$ is the pooled candidate vote share. When sample sizes are in the thousands or tens of thousands, the denominator becomes small. Therefore even several percentage-point differences can produce very large z values. This is why "large vote-share differences in large samples" are difficult to explain by ordinary sampling error alone.

This test has an important assumption: early voters and election-day voters must be random samples from the same voter population. In real elections, this assumption may fail. Early voters may self-select by age, occupation, mobility, party preference, campaign mobilization, and other factors. Therefore a large early-vote/election-day difference alone does not prove misconduct.

Still, the test is not meaningless. Suspicion becomes stronger when:

\begin{enumerate}
\item the early/election-day difference repeats in the same party direction across many constituencies;
\item the difference suddenly becomes larger than in earlier elections in the same country;
\item large residual differences remain after accounting for age, gender, region, urban/rural structure, and early-vote rates;
\item in-district and out-of-district early-vote patterns differ from each other;
\item the same election or region also shows separate anomalies such as identical vote pairs, ballot shortages, or counting-statement inconsistencies.
\end{enumerate}

This paper therefore treats early-vote versus election-day vote-share differences as an auxiliary test. Identical vote-pair repetition is a discrete count-collision problem; early/election-day vote-share difference is a two-sample proportion problem. They have different statistical structures. But if both appear in the same election environment, a single ordinary-randomness explanation becomes weaker.

\subsection{Check Against Official National Assembly Data}

This study directly parses official counting data for the 2016, 2020, and 2024 National Assembly elections. The analyzed constituencies are those in which both a Democratic-party candidate and a conservative-party candidate are present. The conservative party is defined as Saenuri Party in 2016, United Future Party in 2020, and People Power Party in 2024.

For each constituency, the script computes:

\[
\Delta = p_E - p_D
\]

where $p_E$ is the Democratic candidate's share of Democratic-plus-conservative two-party votes in early voting, and $p_D$ is the same two-party share in election-day voting. Early voting combines in-district and out-of-district early votes. Election-day voting includes only election-day polling-station votes. Residential/shipboard, overseas, subtotal, and aggregate rows are excluded.

The output \texttt{outputs/early\_day\_assembly\_summary.csv} reports:

\begin{center}
\scriptsize
\begin{tabular}{l r r r r r r r r}
\hline
Election & Constituencies & Mean diff. & Median & Early adv. & Early dis. & $|z|>5$ & $|z|>10$ & Max $|z|$ \\
\hline
2016 National Assembly & 229 & +3.90\%p & +3.43\%p & 211 & 18 & 158 & 88 & 31.94 \\
2020 National Assembly & 236 & +10.88\%p & +11.21\%p & 236 & 0 & 236 & 234 & 55.77 \\
2024 National Assembly & 245 & +10.35\%p & +10.80\%p & 245 & 0 & 245 & 245 & 60.92 \\
\hline
\end{tabular}
\end{center}

These results show two things at once.

First, the simplified claim that "the law of large numbers held perfectly in 2016 and broke only from 2020" does not fit the official data. In 2016, the Democratic candidate's early-vote two-party share was also higher than the election-day share by an average of 3.90 percentage points, and 211 of 229 constituencies showed a Democratic early-vote advantage.

Second, the 2020 and 2024 patterns are much stronger than 2016. In 2020, all 236 analyzable constituencies showed a Democratic early-vote advantage. In 2024, all 245 analyzable constituencies showed the same direction. The average difference also expanded to roughly 10 to 11 percentage points.

The defensible claim from the direct calculation is therefore:

\begin{quote}
A Democratic early-vote advantage already existed in 2016, but in 2020 and 2024 the direction expanded to every analyzable constituency and the average magnitude grew to about three times the 2016 level.
\end{quote}

This claim is statistically strong. It is still not direct proof of misconduct. Alternative explanations such as voter self-selection, party-specific early-vote mobilization, age, occupation, regional mobility, and COVID-era voting behavior must be evaluated. However, the 2020 and 2024 all-constituency directional pattern is difficult to explain as ordinary sampling error alone, and it is a valid target for independent audit or additional statistical modeling.

The directionality issue is often described as a "law of large numbers" issue. More precisely, the question is whether the same-direction pattern would arise from sampling error if early votes and election-day votes were samples from the same underlying electorate. This paper uses it only as an auxiliary test, with detailed assumptions and results also summarized in Appendix B.

\subsection{Connection to the 2026 Local Election}

The twelve reported identical-vote rows in the June 3, 2026 local election are separate from the early/election-day vote-share difference issue. But both phenomena arise in early-vote data.

This paper focuses on the June 3 local election for the following reasons:

\begin{enumerate}
\item the Songdo 1-dong/Songdo 2-dong identical pair occurred in in-district early voting;
\item the ten Gwangju-Jeonnam rows, or five identical pairs, also occurred in in-district early voting;
\item in past official governor data, the maximum identical-pair repetition within one contest was three pairs;
\item the June 3 reported case reaches five pairs within one contest environment, exceeding the expanded historical benchmark;
\item in the 2020 and 2024 National Assembly elections, early voting also showed systematic directionality relative to election-day voting.
\end{enumerate}

Therefore the June 3 identical-vote-pair problem should not be treated as an isolated coincidence without further checking. Combined with the systematic early-vote deviations observed in 2020 and 2024, it points to a broader audit question: repeated statistical anomalies are appearing in early-vote data, and official raw records are needed to separate lawful voter behavior from administrative or counting-process explanations.

\section{Expected Objections and Responses}

\subsection{"With many counting units, identical vote pairs can always appear."}

Yes, a single duplicate can appear. The paper does not deny that. The question is whether five repeated top-two pairs in the same defined contest environment are ordinary. The historical benchmark and probability calculations show that five is not the same as one.

\subsection{"If 2014 is included, historical duplicates become more common."}

Correct. This is why the paper includes 2014 rather than excluding it. The inclusion makes the benchmark more conservative. Even with 2014 included, the historical maximum within one governor contest is three, not five.

\subsection{"If three historical pairs are acceptable, five should also be acceptable."}

Not under the same probability model. Three or more and five or more have substantially different tail probabilities. Three is a rare but observed historical upper value. Five exceeds that observed upper value and is approximately 37 times rarer than three under the $N=393, K=100,945$ Poisson calculation.

\subsection{"Unit sizes and regional preferences invalidate the simple probability calculation."}

They may affect the calculation, which is why the paper does not rely only on a uniform theoretical model. It estimates $K$ from actual Korean governor-election data and performs resampling from historical actual vote pairs. A critic who argues that regional structure fully explains the anomaly should provide a better model and show that it generates five-pair clusters with non-negligible frequency.

\subsection{"The nationwide twelve cases were found after searching, so the probability is overstated."}

This is a serious objection to the nationwide headline. The paper therefore places the main probability claim on the Gwangju-Jeonnam internal cluster rather than the entire nationwide set. The Songdo case is reported separately.

\subsection{"The 2026 data are not yet from an official integrated file."}

This is the most important limitation. The present package rechecks the event rows from official NEC election-statistics HTML, not from a historical-style integrated XLSX file. Therefore the conclusion is not a final legal determination. It is a statistical audit trigger. When the integrated file or original counting statements are released, the event definition should be reapplied.

\subsection{"Early-vote and election-day differences can be caused by voter self-selection."}

Yes. That is why the early-vote analysis is not used as direct proof of misconduct. It is an auxiliary test. A full explanation should include age, occupation, mobility, region, party mobilization, pandemic-era effects, and other variables.

\subsection{"Rare events do sometimes occur."}

Yes. Statistical rarity alone is not impossibility. But public election administration does not require impossibility before audit. A rare official-result pattern can justify raw-data disclosure, especially when the records needed to resolve the question are held by the election-management institution.

\subsection{"How can the claim be falsified?"}

The claim is falsifiable. It would be weakened if:

1. Official integrated counting-unit files show that the event rows were misreported or misread.
2. Original counting statements and first-pass machine outputs reconcile the final equal pairs through documented review steps.
3. A broader independently reproduced baseline using the same event definition shows that five-pair clusters commonly occur in comparable contests.
4. A covariate-adjusted model explains the early-vote/election-day pattern without leaving systematic residual anomalies.

\section{Raw-Data Audit Decision Criteria}

An audit should not be vague. The minimum decision criteria should be:

\begin{itemize}
\item Do the official integrated files match the HTML page values?
\item Do original counting statements match the published results?
\item Do ballot-sorter first-pass results differ from final results, and if so, are all differences explained by review logs?
\item Are reviewed ballots allocated by candidate in a way that reproduces the final equality?
\item Are data-entry and correction logs consistent with the published timeline?
\item Are all twelve event rows traceable from ballot-level or batch-level records to final publication?
\end{itemize}

If these criteria are met, the anomaly can be administratively explained. If they are not met, the anomaly remains unresolved.

\section{Conclusion}

The 2026 Gwangju-Jeonnam repeated identical in-district early-vote pairs are not adequately explained by ordinary chance under the tests reported here. The historical governor benchmark contains 51 constituency contests and 1,514,172 actual top-two comparisons, yet no contest reaches five repeated pairs. The Poisson approximation gives about 0.115\% for five or more repeated pairs under $N=393$ and $K=100,944.8$. The nonparametric resampling test produces zero five-pair trials in 200,000 repetitions.

The Incheon Songdo case adds a separate low-probability exact match whose final equality appears alongside different auxiliary counts. The early-vote/election-day analysis adds a broader directional pattern in recent National Assembly elections. None of these facts alone identifies a culprit or proves a criminal act. Together, however, they create a statistically serious anomaly that cannot be closed by the word "coincidence" without raw evidence.

The proper conclusion is therefore: the current public data do not prove a legal finding of misconduct, but they do establish a strong statistical basis for official raw-data disclosure and independent audit.

\section{Data and Code Availability}

The reproduction package contains the historical public election files, NEC HTML cache for the 2026 checks, parsing scripts, probability scripts, output CSVs, checksums, the Korean paper, this English paper, and compiled PDFs. The main reproduction command is:

\begin{verbatim}
python3 scripts/run_all.py
\end{verbatim}

The submission package can be rebuilt with:

\begin{verbatim}
python3 scripts/create_submission_zip.py
\end{verbatim}

Package validation can be run with:

\begin{verbatim}
python3 scripts/validate_package.py
\end{verbatim}

\section{Research Ethics and Conflict of Interest}

This study uses public election data and public reports. It does not process private voter-level data. The author declares no institutional affiliation. The purpose of the study is to define a reproducible statistical anomaly and the raw records needed to resolve it. The paper should not be read as a legal accusation against a specific person without the additional raw administrative records identified above.

\section*{Appendix A. Intuitive Meaning of Key Terms}

\begin{center}
\small
\begin{longtable}{p{0.23\textwidth} p{0.65\textwidth}}
\hline
Term & Intuitive meaning \\
\hline
Null hypothesis & The default assumption that the observed pattern could arise by chance \\
Alternative hypothesis & The competing assumption that the pattern is difficult to treat as ordinary chance and requires another explanation \\
Poisson approximation & A standard approximation used to calculate how often rare events repeat \\
Counting unit & The smallest official result row used for comparison, such as an eup, myeon, or dong in-district early-vote row \\
In-district early voting & Early votes cast by voters within their registered district \\
Identical vote pair & Two counting units have the same first-candidate count and same second-candidate count \\
Effective pair space $K$ & The practical range of first- and second-place candidate vote-count combinations that can actually appear in in-district early voting \\
Sensitivity analysis & A check of whether the conclusion remains similar when assumptions such as $K$ or $N$ are changed \\
Nonparametric resampling without replacement & A simulation that draws actual historical counting units without drawing the same unit twice \\
Rule of three & A conservative rule of thumb: if an event is observed zero times in $n$ trials, its rough 95\% upper bound is $3/n$ \\
Tail probability & The probability of observing an extreme event far from the average case \\
Law of large numbers & The principle that, as sample size grows, the proportion in a random sample approaches the population proportion \\
Sign test & A nonparametric test that focuses on how often the direction of a difference repeats, rather than the size of the difference \\
Audit trigger & A signal strong enough to require raw-data disclosure and independent verification, without itself proving misconduct \\
\hline
\end{longtable}
\end{center}

\section*{Appendix B. Early Voting and the Law of Large Numbers}

Public discussions often invoke the "law of large numbers" in early-vote debates, but that phrase alone cannot prove misconduct. The law of large numbers says that as sample size grows, the proportion in a random sample tends to approach the population proportion. To apply that principle to early voting, early voters and election-day voters would have to be randomly split from the same population. In real elections, this premise is incomplete because voters choose whether to vote early.

Nevertheless, the law-of-large-numbers discussion is not meaningless. If sample sizes are large and nearly every constituency shows an early-vote advantage in the same party direction, that pattern is difficult to treat as sampling error alone. This appendix checks that point using the simplest sign test.

For each constituency, the calculation asks whether the Democratic candidate's early-vote two-party share is higher than the same candidate's election-day two-party share. The null hypothesis is:

\begin{quote}
If the early/election-day difference is not tilted toward one party direction, the probability that a constituency shows a Democratic early-vote advantage is roughly $1/2$.
\end{quote}

Under this null hypothesis, the probability that at least $x$ out of $n$ constituencies show a Democratic early-vote advantage is the one-sided tail probability:

\[
P(X \ge x) = \sum_{k=x}^{n} \binom{n}{k}\left(\frac{1}{2}\right)^n
\]

The same idea can be described in terms of the random variable:

\[
\Delta = p_E - p_D
\]

where $\Delta$ is the difference between the early-vote share and election-day share. If simple sampling error is the main cause, the distribution of $\Delta$ should fluctuate around zero. If most constituencies have $\Delta>0$, and the center of the distribution moves toward roughly +10 percentage points, the pattern indicates systematic directionality. Whether that directionality is caused by voter self-selection, party mobilization, administrative processes, or another mechanism requires separate data. But the claim that it is not a statistical issue is difficult to sustain.

Direct parsing of official National Assembly data gives:

\begin{center}
\scriptsize
\begin{tabular}{l r r r r r}
\hline
Election & Constituencies & Early adv. & Early dis. & One-sided sign-test p-value & Intuitive frequency \\
\hline
2016 National Assembly & 229 & 211 & 18 & $2.99 \times 10^{-43}$ & about 1 in $3.35 \times 10^{42}$ \\
2020 National Assembly & 236 & 236 & 0 & $9.06 \times 10^{-72}$ & about 1 in $1.10 \times 10^{71}$ \\
2024 National Assembly & 245 & 245 & 0 & $1.77 \times 10^{-74}$ & about 1 in $5.65 \times 10^{73}$ \\
\hline
\end{tabular}
\end{center}

The meaning of this calculation is limited but clear. If early/election-day differences were merely moving up and down like sampling error, the probability of seeing Democratic early-vote advantages in all 236 analyzable constituencies in 2020 and all 245 analyzable constituencies in 2024 would be extremely small under this simple model. Therefore "it only looked that way because the sample was large" is not a sufficient explanation under the symmetric sign-test null.

This result must not be read as direct proof of misconduct. The sign test does not remove alternative explanations such as voter self-selection, party-specific early-vote mobilization, age, occupation, regional mobility, or COVID-period voting effects. If those factors were strongly operating, the symmetric $1/2$ null hypothesis would be weakened.

The appendix conclusion is therefore narrow. The same-direction early-vote pattern in the 2020 and 2024 National Assembly elections is statistically extremely unlikely under a simple sampling-error model. But it is an auxiliary anomaly indicator that supports the identical-vote-pair analysis; legal or administrative conclusions require constituency-level voter composition, campaign-mobilization evidence, separated in-district/out-of-district early-vote data, original counting statements, and first-pass sorter records.

\section*{References}

\begin{itemize}
\item National Election Commission, election-statistics system, counting-unit result pages.
\item Public Data Portal, National Election Commission vote-counting data files for past elections.
\item Public Data Portal, 2014 nationwide local election counting data.
\item Public Data Portal, 2016 National Assembly election counting data.
\item Public Data Portal, 2018 nationwide local election counting data.
\item Public Data Portal, 2020 National Assembly election counting data.
\item Public Data Portal, 2022 nationwide local election counting data.
\item Public Data Portal, 2024 National Assembly election counting data.
\item News1/Daum, June 6, 2026, report on identical vote counts in Songdo 1-dong and Songdo 2-dong.
\item Hankyoreh, June 8, 2026, report on the Songdo identical-count controversy and election-management explanation.
\item TV Chosun, June 8, 2026, report on identical in-district early-vote counts in twelve places nationwide.
\end{itemize}

\section*{Reproduction Outputs}

Key files in the package include:

\begin{itemize}
\item \texttt{outputs/dataset\_counts.csv}
\item \texttt{outputs/duplicate\_summary.csv}
\item \texttt{outputs/governor\_actual\_top2\_summary.csv}
\item \texttt{outputs/governor\_actual\_top2\_by\_contest.csv}
\item \texttt{outputs/governor\_actual\_top2\_duplicates.csv}
\item \texttt{outputs/governor\_bootstrap\_summary.csv}
\item \texttt{outputs/governor\_bootstrap\_histogram.csv}
\item \texttt{outputs/nec\_2026\_gwangju\_jeonnam\_unit\_counts.csv}
\item \texttt{outputs/nec\_2026\_gwangju\_jeonnam\_units.csv}
\item \texttt{outputs/nec\_2026\_gwangju\_jeonnam\_unit\_summary.json}
\item \texttt{outputs/nec\_2026\_reported\_duplicate\_cases.csv}
\item \texttt{outputs/nec\_2026\_reported\_duplicate\_pairs.csv}
\item \texttt{outputs/songdo\_probability\_summary.csv}
\item \texttt{outputs/songdo\_official\_rows.csv}
\item \texttt{outputs/probability\_core.csv}
\item \texttt{outputs/probability\_k\_sensitivity.csv}
\item \texttt{outputs/probability\_n\_sensitivity.csv}
\item \texttt{outputs/early\_day\_assembly\_summary.csv}
\item \texttt{outputs/early\_day\_assembly\_twoparty.csv}
\item \texttt{outputs/checksums\_sha256.csv}
\end{itemize}

\end{document}
